\section{Описание программы}

\subsection{Архитектура проекта}

Программа реализована на основе класса \texttt{Polynomial}, обеспечивающего работу с многочленом через динамический массив коэффициентов \texttt{double*}. Коэффициенты хранятся по степеням: \texttt{A[0]} — свободный член, \texttt{A[i]} — коэффициент при $x^i$. Управление памятью осуществляется вручную через \texttt{new[]}/\texttt{delete[]}, соблюдается RAII.

\subsection{Состав класса}

Класс \texttt{Polynomial} содержит:

\begin{table}[htbp]
\centering
\caption{Методы класса \texttt{Polynomial}}
\begin{tabular}{@{}>{\raggedright\arraybackslash}p{0.96\textwidth}@{}}
\hline
\textbf{Сигнатура} \\
\hline
\texttt{Polynomial(const double* data, int degree)} \\
\texttt{Polynomial(const Polynomial\& other)} \\
\texttt{Polynomial(Polynomial\&\& other) noexcept} \\
\texttt{Polynomial\& operator=(Polynomial other) noexcept} \\
\texttt{\~{}Polynomial()} \\
\texttt{void print(std::ostream\& os) const} \\
\texttt{double evaluate(double in) const} \\
\texttt{void evaluate\_many(const double* xs, double* ys, size\_t n) const} \\
\texttt{void add(const Polynomial\& other)} \\
\texttt{void mul(const Polynomial\& other)} \\
\hline
\end{tabular}
\end{table}

Вспомогательные методы: \texttt{void trim()} (удаление старших нулевых коэффициентов), \texttt{void swap(Polynomial\& other) noexcept}.

\subsection{Особенности реализации}

\begin{itemize}
  \item Старшие нулевые коэффициенты отбрасываются методом \texttt{trim()} после операций, что исключает хранение лишних членов.
  \item При всех коэффициентах, равных нулю, многочлен считается пустым (\texttt{degree\_ = -1}, \texttt{data\_ = nullptr}).
  \item Оператор присваивания реализован через copy-and-swap.
  \item Конструктор копирования выполняет глубокое копирование.
  \item Перемещающий конструктор передаёт владение памятью без копирования.
\end{itemize}

\subsection{Сборка и запуск}

Проект собирается out-of-source через CMake: статическая библиотека \texttt{polynomial}, исполняемый файл \texttt{main} и тестовый исполняемый файл \texttt{test\_polynomial}. Google Test подключается через \texttt{FetchContent} (версия 1.14.0), тесты регистрируются в CTest командой \texttt{gtest\_discover\_tests}.

\begin{lstlisting}[style=cmakestyle,caption={Сборка и запуск},numbers=none]
cmake -S . -B build
cmake --build build
./build/main
ctest --test-dir build --output-on-failure
\end{lstlisting}

\subsection{Исходный код}

\subsubsection{Система сборки}

\lstinputlisting[style=cmakestyle,caption={CMakeLists.txt}]{../../CMakeLists.txt}

\subsubsection{Заголовочный файл}

\lstinputlisting[style=cppstyle,caption={src/polynomial.h}]{../../src/polynomial.h}

\subsubsection{Реализация класса}

\lstinputlisting[style=cppstyle,caption={src/polynomial.cpp}]{../../src/polynomial.cpp}

\subsubsection{Точка входа}

\lstinputlisting[style=cppstyle,caption={src/main.cpp}]{../../src/main.cpp}

\subsubsection{Тесты}

\lstinputlisting[style=cppstyle,caption={tests/test\_polynomial.cpp}]{../../tests/test_polynomial.cpp}
